% Options for packages loaded elsewhere
\PassOptionsToPackage{unicode}{hyperref}
\PassOptionsToPackage{hyphens}{url}
\documentclass[
]{article}
\usepackage{xcolor}
\usepackage{amsmath,amssymb}
\setcounter{secnumdepth}{5}
\usepackage{iftex}
\ifPDFTeX
  \usepackage[T1]{fontenc}
  \usepackage[utf8]{inputenc}
  \usepackage{textcomp} % provide euro and other symbols
\else % if luatex or xetex
  \usepackage{unicode-math} % this also loads fontspec
  \defaultfontfeatures{Scale=MatchLowercase}
  \defaultfontfeatures[\rmfamily]{Ligatures=TeX,Scale=1}
\fi
% \usepackage{lmodern}
\ifPDFTeX\else
  % xetex/luatex font selection
\fi
% Use upquote if available, for straight quotes in verbatim environments
\IfFileExists{upquote.sty}{\usepackage{upquote}}{}
\IfFileExists{microtype.sty}{% use microtype if available
  \usepackage[]{microtype}
  \UseMicrotypeSet[protrusion]{basicmath} % disable protrusion for tt fonts
}{}
\makeatletter
\@ifundefined{KOMAClassName}{% if non-KOMA class
  \IfFileExists{parskip.sty}{%
    \usepackage{parskip}
  }{% else
    \setlength{\parindent}{0pt}
    \setlength{\parskip}{6pt plus 2pt minus 1pt}}
}{% if KOMA class
  \KOMAoptions{parskip=half}}
\makeatother
\usepackage{color}
\usepackage{fancyvrb}
\newcommand{\VerbBar}{|}
\newcommand{\VERB}{\Verb[commandchars=\\\{\}]}
\DefineVerbatimEnvironment{Highlighting}{Verbatim}{commandchars=\\\{\}}
% Add ',fontsize=\small' for more characters per line
\newenvironment{Shaded}{}{}
\newcommand{\AlertTok}[1]{\textcolor[rgb]{1.00,0.00,0.00}{\textbf{#1}}}
\newcommand{\AnnotationTok}[1]{\textcolor[rgb]{0.38,0.63,0.69}{\textbf{\textit{#1}}}}
\newcommand{\AttributeTok}[1]{\textcolor[rgb]{0.49,0.56,0.16}{#1}}
\newcommand{\BaseNTok}[1]{\textcolor[rgb]{0.25,0.63,0.44}{#1}}
\newcommand{\BuiltInTok}[1]{\textcolor[rgb]{0.00,0.50,0.00}{#1}}
\newcommand{\CharTok}[1]{\textcolor[rgb]{0.25,0.44,0.63}{#1}}
\newcommand{\CommentTok}[1]{\textcolor[rgb]{0.38,0.63,0.69}{\textit{#1}}}
\newcommand{\CommentVarTok}[1]{\textcolor[rgb]{0.38,0.63,0.69}{\textbf{\textit{#1}}}}
\newcommand{\ConstantTok}[1]{\textcolor[rgb]{0.53,0.00,0.00}{#1}}
\newcommand{\ControlFlowTok}[1]{\textcolor[rgb]{0.00,0.44,0.13}{\textbf{#1}}}
\newcommand{\DataTypeTok}[1]{\textcolor[rgb]{0.56,0.13,0.00}{#1}}
\newcommand{\DecValTok}[1]{\textcolor[rgb]{0.25,0.63,0.44}{#1}}
\newcommand{\DocumentationTok}[1]{\textcolor[rgb]{0.73,0.13,0.13}{\textit{#1}}}
\newcommand{\ErrorTok}[1]{\textcolor[rgb]{1.00,0.00,0.00}{\textbf{#1}}}
\newcommand{\ExtensionTok}[1]{#1}
\newcommand{\FloatTok}[1]{\textcolor[rgb]{0.25,0.63,0.44}{#1}}
\newcommand{\FunctionTok}[1]{\textcolor[rgb]{0.02,0.16,0.49}{#1}}
\newcommand{\ImportTok}[1]{\textcolor[rgb]{0.00,0.50,0.00}{\textbf{#1}}}
\newcommand{\InformationTok}[1]{\textcolor[rgb]{0.38,0.63,0.69}{\textbf{\textit{#1}}}}
\newcommand{\KeywordTok}[1]{\textcolor[rgb]{0.00,0.44,0.13}{\textbf{#1}}}
\newcommand{\NormalTok}[1]{#1}
\newcommand{\OperatorTok}[1]{\textcolor[rgb]{0.40,0.40,0.40}{#1}}
\newcommand{\OtherTok}[1]{\textcolor[rgb]{0.00,0.44,0.13}{#1}}
\newcommand{\PreprocessorTok}[1]{\textcolor[rgb]{0.74,0.48,0.00}{#1}}
\newcommand{\RegionMarkerTok}[1]{#1}
\newcommand{\SpecialCharTok}[1]{\textcolor[rgb]{0.25,0.44,0.63}{#1}}
\newcommand{\SpecialStringTok}[1]{\textcolor[rgb]{0.73,0.40,0.53}{#1}}
\newcommand{\StringTok}[1]{\textcolor[rgb]{0.25,0.44,0.63}{#1}}
\newcommand{\VariableTok}[1]{\textcolor[rgb]{0.10,0.09,0.49}{#1}}
\newcommand{\VerbatimStringTok}[1]{\textcolor[rgb]{0.25,0.44,0.63}{#1}}
\newcommand{\WarningTok}[1]{\textcolor[rgb]{0.38,0.63,0.69}{\textbf{\textit{#1}}}}
\usepackage{longtable,booktabs,array}
\newcounter{none} % for unnumbered tables
\usepackage{calc} % for calculating minipage widths
% Correct order of tables after \paragraph or \subparagraph
\usepackage{etoolbox}
\makeatletter
\patchcmd\longtable{\par}{\if@noskipsec\mbox{}\fi\par}{}{}
\makeatother
% Allow footnotes in longtable head/foot
\IfFileExists{footnotehyper.sty}{\usepackage{footnotehyper}}{\usepackage{footnote}}
\makesavenoteenv{longtable}
\usepackage{graphicx}
\makeatletter
\newsavebox\pandoc@box
\newcommand*\pandocbounded[1]{% scales image to fit in text height/width
  \sbox\pandoc@box{#1}%
  \Gscale@div\@tempa{\textheight}{\dimexpr\ht\pandoc@box+\dp\pandoc@box\relax}%
  \Gscale@div\@tempb{\linewidth}{\wd\pandoc@box}%
  \ifdim\@tempb\p@<\@tempa\p@\let\@tempa\@tempb\fi% select the smaller of both
  \ifdim\@tempa\p@<\p@\scalebox{\@tempa}{\usebox\pandoc@box}%
  \else\usebox{\pandoc@box}%
  \fi%
}
% Set default figure placement to htbp
\def\fps@figure{htbp}
\makeatother
\setlength{\emergencystretch}{3em} % prevent overfull lines
\providecommand{\tightlist}{%
  \setlength{\itemsep}{0pt}\setlength{\parskip}{0pt}}
\usepackage[]{natbib}
\bibliographystyle{plainnat}
\makeatletter
\@ifpackageloaded{subcaption}{}{\usepackage{subcaption}}
\@ifpackageloaded{caption}{}{\usepackage{caption}}
\captionsetup[subfigure]{margin=0.5em}
\AtBeginDocument{%
\renewcommand*\figurename{Figure}
\renewcommand*\tablename{Table}
}
\AtBeginDocument{%
\renewcommand*\listfigurename{List of Figures}
\renewcommand*\listtablename{List of Tables}
}
\newcounter{pandoccrossref@subfigures@footnote@counter}
\newenvironment{pandoccrossrefsubfigures}{%
\setcounter{pandoccrossref@subfigures@footnote@counter}{0}
\begin{figure}\centering%
\gdef\global@pandoccrossref@subfigures@footnotes{}%
\DeclareRobustCommand{\footnote}[1]{\footnotemark%
\stepcounter{pandoccrossref@subfigures@footnote@counter}%
\ifx\global@pandoccrossref@subfigures@footnotes\empty%
\gdef\global@pandoccrossref@subfigures@footnotes{{##1}}%
\else%
\g@addto@macro\global@pandoccrossref@subfigures@footnotes{, {##1}}%
\fi}}%
{\end{figure}%
\addtocounter{footnote}{-\value{pandoccrossref@subfigures@footnote@counter}}
\@for\f:=\global@pandoccrossref@subfigures@footnotes\do{\stepcounter{footnote}\footnotetext{\f}}%
\gdef\global@pandoccrossref@subfigures@footnotes{}}
\@ifpackageloaded{float}{}{\usepackage{float}}
\floatstyle{ruled}
\@ifundefined{c@chapter}{\newfloat{codelisting}{h}{lop}}{\newfloat{codelisting}{h}{lop}[chapter]}
\floatname{codelisting}{Listing}
\newcommand*\listoflistings{\listof{codelisting}{List of Listings}}
\makeatother
\usepackage{bookmark}
\IfFileExists{xurl.sty}{\usepackage{xurl}}{} % add URL line breaks if available
\urlstyle{same}
% fallback for those not using the hyperref driver hyperxmp:
\makeatletter
\@ifundefined{xmpquote}{\newcommand{\xmpquote}[1]{#1}}{}
\makeatother
\hypersetup{
  pdftitle={Data Descriptor Template: Schema{,} Provenance{,} and Release Readiness},
  pdfauthor={Daniel Ari Friedman},
  colorlinks=true,linkcolor=red,urlcolor=red,citecolor=red,anchorcolor=red,filecolor=red,
  pdfcreator={LaTeX via pandoc}}

\title{Data Descriptor Template: Schema, Provenance, and Release
Readiness}
\author{Daniel Ari Friedman}
\date{}

\IfFileExists{seqsplit.sty}{\usepackage{seqsplit}}{\newcommand{\seqsplit}[1]{#1}}
\protected\def\breakseq#1{\seqsplit{#1}}
\protected\def\breaktt#1{\begingroup\ttfamily\seqsplit{#1}\endgroup}
\title{Data Descriptor Template: Schema, Provenance, and Release Readiness}
\author{Daniel Ari Friedman\\\footnotesize{Active Inference Institute}\\\footnotesize{\texttt{daniel@activeinference.institute}}\\\footnotesize{\href{https://orcid.org/0000-0001-6232-9096}{ORCID: 0000-0001-6232-9096}} \\ \footnotesize{DOI: \href{https://doi.org/10.5281/zenodo.21298883}{10.5281/zenodo.21298883}}}
\date{2026-07-10}

\usepackage[margin=0.75in]{geometry}
\usepackage{graphicx}
\setlength{\parindent}{0pt}
\setlength{\parskip}{0.45em}

\begin{document}
\begin{titlepage}
\centering
\vspace*{2cm}
{\LARGE\bfseries Data Descriptor Template: Schema, Provenance, and Release Readiness\par}
\vspace{0.75em}
{\large A FAIR-style exemplar for dataset papers\par}
\vfill
\makeatletter
{\@author\par}
\vspace{1em}
{\@date\par}
\makeatother
\vfill
\end{titlepage}
\thispagestyle{empty}
\tableofcontents
\newpage

\section{Abstract}\label{abstract}

This exemplar demonstrates a \textbf{data descriptor} workflow in which
the schema, file inventory, provenance chain, license boundary, and
validation gate are treated as first-class research artifacts rather
than afterthoughts. It ships a small, public, synthetic demonstration
dataset (two CSV files under \texttt{data/fixtures/}) and a
machine-readable descriptor (\breaktt{data/example\_descriptor.json})
that declares each file's media type, sha256 checksum, and row count
alongside a six-field data dictionary with typed constraints. A tested
validation library (\breaktt{src/data\_descriptor/}) checks the
descriptor's shape, safety, and completeness; recomputes each declared
checksum and row count against the bytes on disk; and emits a
deterministic, metadata-only release manifest suitable for
pre-publication review. Every figure and quantitative claim in this
manuscript is produced by that library and regenerated on demand, so the
prose describes structure and provenance rather than transcribing values
that would drift. This is a template with a demonstration dataset: it
makes no scientific claim about the data, only about how to describe and
release a dataset responsibly.

\newpage

\section{Introduction}\label{sec:introduction}

A \emph{data descriptor} (or data paper) publishes a dataset as a
citable research object. Unlike an analysis paper, its contribution is
not a statistical finding but a \textbf{contract}: a precise,
machine-readable account of what the dataset contains, how it was
produced, how it is licensed, and what quality guarantees it carries.
The FAIR guiding principles --- Findable, Accessible, Interoperable,
Reusable \citep{wilkinson2016fair} --- frame why this matters: a dataset
is only reusable to the extent that its structure and provenance are
legible to both humans and machines.

\subsection{What this template
provides}\label{what-this-template-provides}

This exemplar packages the recurring moving parts of a data descriptor
so that a fork can start from a working, tested baseline:

\begin{itemize}
\tightlist
\item
  A \textbf{schema / data dictionary}: named fields with declared types,
  nullability, units, and value constraints (patterns, enumerations,
  numeric bounds).
\item
  A \textbf{file inventory}: each shipped file with its media type, a
  sha256 checksum, and a row count.
\item
  A \textbf{provenance chain}: the ordered steps and agents that
  produced the release.
\item
  A \textbf{license boundary}: an explicit reuse license, and the
  discipline of publishing \emph{metadata and checksums} rather than
  bundling restricted bytes.
\item
  A \textbf{validation gate}: a tested library that rejects malformed,
  unsafe, or incomplete descriptors and verifies declared checksums
  against real bytes.
\end{itemize}

\subsection{Scope and honesty}\label{scope-and-honesty}

The shipped dataset is \textbf{synthetic and deliberately small} --- it
exists to make the workflow runnable and testable, not to support any
empirical claim. Its values are generated deterministically.
Accordingly, this manuscript restricts its claims to dataset
\emph{structure, provenance, quality, and release readiness}. When a
real dataset is forked in, the same validation and figure code applies
unchanged; only the descriptor and fixture files are replaced. The
sections that follow describe the demonstration dataset
(sec.~\ref{sec:data}), its schema (sec.~\ref{sec:schema}), its
provenance (sec.~\ref{sec:provenance}), and the quality gate that binds
the descriptor to the bytes it names (sec.~\ref{sec:quality}), and close
with usage and forking notes (sec.~\ref{sec:usage}).

\newpage

\section{Data description}\label{sec:data}

The demonstration dataset consists of two comma-separated files under
\texttt{data/fixtures/}, described by the machine-readable descriptor
\breaktt{data/example\_descriptor.json}:

\begin{itemize}
\tightlist
\item
  \breaktt{fixtures/measurements.csv} --- one row per synthetic sample,
  with a bounded measurement, an assignment group, a capture instrument,
  and a collection date.
\item
  \breaktt{fixtures/subjects.csv} --- one row per synthetic subject, with
  an enrollment site and date. Each measurement references a subject
  through the \texttt{subject\_id} foreign key.
\end{itemize}

\subsection{File inventory}\label{file-inventory}

The descriptor declares each file's media type, sha256 checksum, and row
count. fig.~\ref{fig:file_inventory} shows the declared row counts, read
directly from the descriptor by \breaktt{file\_inventory\_rows()} in
\breaktt{src/data\_descriptor/figures.py} and plotted by the thin
analysis script
\href{../scripts/generate_figures.py}{\breaktt{scripts/generate\_figures.py}}.

\begin{figure}
\centering
\pandocbounded{\includegraphics[keepaspectratio,alt={File inventory: declared row counts per file, from file\_inventory\_rows(). The measurement table is the larger of the two files; the subject table is its dimension companion. Both are declared as text/csv. The bar values are the row counts recorded in the descriptor, not re-derived here --- the quality gate (sec.~) is where declared and actual counts are reconciled.}]{../figures/file_inventory.png}}
\caption{File inventory: declared row counts per file, from
\protect\breaktt{file\_inventory\_rows()}. The measurement table is the
larger of the two files; the subject table is its dimension companion.
Both are declared as \protect\texttt{text/csv}. The bar values are the
row counts recorded in the descriptor, not re-derived here --- the
quality gate (sec.~\ref{sec:quality}) is where declared and actual
counts are reconciled.}\label{fig:file_inventory}
\end{figure}

Running the script regenerates the figures under
\breaktt{manuscript/figures/}; the prose below and in later sections
describes what those artifacts show. Because the figures are produced
from the descriptor and the fixture bytes, they cannot drift from the
data without the tests and the quality gate noticing.

\subsection{Release boundary}\label{release-boundary}

The descriptor is deliberately a \emph{metadata} object. For a real
dataset, the descriptor and its checksums can be published even when the
underlying bytes are access-controlled: the checksums let a downstream
user verify integrity once they obtain the files through the appropriate
channel. In this template the fixture bytes are public and small, so
they are shipped alongside the descriptor and used to demonstrate
end-to-end verification.

\newpage

\section{Schema and data dictionary}\label{sec:schema}

The descriptor's \texttt{fields} list is the dataset's \textbf{data
dictionary}: it names each column of the primary measurement table and
declares its type, nullability, optional unit, and value constraints.
fig.~\ref{fig:schema_overview} renders this dictionary directly from the
descriptor via \breaktt{schema\_table\_rows()} in
\breaktt{src/data\_descriptor/figures.py}.

\begin{figure}
\centering
\pandocbounded{\includegraphics[keepaspectratio,alt={Field schema / data dictionary: one row per declared field, with type, nullability, unit, and a compact constraint label, from schema\_table\_rows(). Identifier fields carry regular-expression patterns; the categorical fields carry closed enumerations; the quantitative field carries a unit and numeric bounds. The table is generated from data/example\_descriptor.json and cannot silently disagree with it.}]{../figures/schema_overview.png}}
\caption{Field schema / data dictionary: one row per declared field,
with type, nullability, unit, and a compact constraint label, from
\protect\breaktt{schema\_table\_rows()}. Identifier fields carry
regular-expression patterns; the categorical fields carry closed
enumerations; the quantitative field carries a unit and numeric bounds.
The table is generated from
\protect\breaktt{data/example\_descriptor.json} and cannot silently
disagree with it.}\label{fig:schema_overview}
\end{figure}

\subsection{Field contract}\label{field-contract}

The six declared fields exercise the constraint vocabulary the validator
understands:

\begin{itemize}
\tightlist
\item
  \textbf{\texttt{sample\_id}} (\texttt{string}, not null) --- the
  primary key, constrained to the pattern
  \texttt{\^{}S{[}0-9{]}\{3\}\$}.
\item
  \textbf{\texttt{subject\_id}} (\texttt{string}, not null) --- foreign
  key into \breaktt{fixtures/subjects.csv}, constrained to
  \texttt{\^{}P{[}0-9{]}\{3\}\$}.
\item
  \textbf{\texttt{group}} (\texttt{category}, not null) --- a closed
  enumeration of \texttt{control} and \texttt{treatment}.
\item
  \textbf{\texttt{value}} (\texttt{number}, not null) --- the
  measurement, carrying the unit \breaktt{normalized\_score} and numeric
  bounds \texttt{{[}0,\ 1{]}}.
\item
  \textbf{\texttt{collected\_on}} (\texttt{date}, not null) --- an
  ISO-8601 collection date.
\item
  \textbf{\texttt{instrument}} (\texttt{category}, not null) --- a
  closed enumeration of \texttt{sensor\_a} and \texttt{sensor\_b}.
\end{itemize}

\subsection{Schema fingerprint}\label{schema-fingerprint}

\breaktt{descriptor\_fingerprint()} reduces the field list to a stable,
order-independent sha256 fingerprint over
\breaktt{(name,\ type,\ nullable)} triples. Reordering the fields leaves
the fingerprint unchanged, so two descriptors with the same schema hash
identically regardless of authoring order --- the property is asserted
directly in the test suite. The fingerprint is what lets a release
manifest reference ``the schema'' by a single value that changes if and
only if the field contract changes.

\subsection{Constraint validation}\label{constraint-validation}

The validator flags quantitative fields that lack a unit, categorical
fields that lack allowed values, enumerations that are present but
empty, numeric bounds where the minimum exceeds the maximum, and
patterns declared on non-text fields. On the shipped descriptor none of
these fire --- the data dictionary above is complete, so the readiness
score is unpenalised (see sec.~\ref{sec:quality}).

\newpage

\section{Provenance}\label{sec:provenance}

Provenance records \emph{how the release came to be}. The descriptor's
\texttt{provenance} list is an ordered chain of steps, each naming the
agent responsible. fig.~\ref{fig:provenance_flow} renders that chain via
\breaktt{provenance\_steps()} in
\breaktt{src/data\_descriptor/figures.py}.

\begin{figure}
\centering
\pandocbounded{\includegraphics[keepaspectratio,alt={Provenance chain: the ordered steps that produced the release, from provenance\_steps(). Each box is a declared step; the label beneath is the responsible agent. The chain runs left to right from raw collection to the packaged, metadata-only release manifest.}]{../figures/provenance_flow.png}}
\caption{Provenance chain: the ordered steps that produced the release,
from \protect\breaktt{provenance\_steps()}. Each box is a declared step;
the label beneath is the responsible agent. The chain runs left to right
from raw collection to the packaged, metadata-only release
manifest.}\label{fig:provenance_flow}
\end{figure}

\subsection{Methods: the provenance
chain}\label{methods-the-provenance-chain}

The provenance chain is this descriptor's methods section --- it records
the generation and processing methodology step by step. The shipped
descriptor declares four steps:

\begin{enumerate}
\def\labelenumi{\arabic{enumi}.}
\tightlist
\item
  \textbf{collect} --- a synthetic fixture generator emits the
  deterministic measurement and subject tables.
\item
  \textbf{clean} --- identifiers are normalized and the measurement
  range is bounded.
\item
  \textbf{validate} --- the descriptor is checked for schema,
  constraint, and byte-level agreement.
\item
  \textbf{package} --- the metadata-only release manifest is emitted.
\end{enumerate}

\subsection{Why depth matters}\label{why-depth-matters}

The validator treats provenance shorter than two steps as a warning ---
a release that records only ``we have the data'' but not ``how it was
produced and checked'' is thinner than a reusable descriptor should be.
The four-step chain here pairs a collection origin with an explicit
validation and packaging trail, which is the minimum a downstream reuser
needs to judge whether the dataset was produced the way they require.
For a real dataset, these steps would cite concrete tools, versions, and
operators rather than the template's synthetic agents.

\newpage

\section{Quality validation}\label{sec:quality}

The quality gate is what separates a descriptor that merely
\emph{claims} to describe a dataset from one that \emph{provably} does.
It has two layers: structural validation of the descriptor, and
byte-level reconciliation of the descriptor against the files it names.

\subsection{The validation gate reacts to
defects}\label{the-validation-gate-reacts-to-defects}

\breaktt{validate\_descriptor()} emits severity-tagged findings for
missing required keys, unsafe or duplicate file paths, malformed
checksums, unknown media types, duplicate or malformed fields,
constraint violations, missing primary keys, and thin provenance.
\breaktt{build\_descriptor\_report()} folds these into a readiness score,
penalising each error and warning.

To show the gate is not vacuous, fig.~\ref{fig:quality_gate} compares
the shipped fixture against a deliberately-perturbed copy produced by
\breaktt{demo\_broken\_descriptor()} (which strips the license, corrupts
a checksum, zeroes a row count, and removes a unit). The clean fixture
produces zero findings; the perturbed demo produces several errors and
warnings.

\begin{figure}
\centering
\pandocbounded{\includegraphics[keepaspectratio,alt={Quality gate: validation findings by severity for the clean fixture descriptor (left, zero findings) versus a deliberately-broken demonstration copy (right). The demonstration perturbation is clearly named in code and never treated as real data --- it exists only to prove the gate reacts. Counts come from severity\_counts().}]{../figures/quality_gate.png}}
\caption{Quality gate: validation findings by severity for the clean
fixture descriptor (left, zero findings) versus a deliberately-broken
demonstration copy (right). The demonstration perturbation is clearly
named in code and never treated as real data --- it exists only to prove
the gate reacts. Counts come from
\protect\breaktt{severity\_counts()}.}\label{fig:quality_gate}
\end{figure}

\subsection{Descriptor\texorpdfstring{\ensuremath{\leftrightarrow}}{<->}file
verification}\label{descriptorfile-verification}

A checksum is only meaningful if something checks it.
\breaktt{verify\_descriptor\_files()} recomputes the sha256 digest and
data-row count of every declared file that is present on disk and
compares them to the descriptor's declarations, reporting each file as
\texttt{verified}, \breaktt{checksum\_mismatch}, \texttt{row\_mismatch},
or \texttt{absent}. Files a descriptor references but does not bundle
are reported as \texttt{absent} rather than as failures --- verification
is asserted only for bytes that are actually present, which is what
makes the metadata-only release boundary honest.

fig.~\ref{fig:checksum_verification} shows this reconciliation for the
shipped fixture: both files verify, declared and actual row counts
agree, and each recomputed checksum matches its declaration, so the
readiness score is unpenalised.

\begin{figure}
\centering
\pandocbounded{\includegraphics[keepaspectratio,alt={Descriptor\texorpdfstring{\ensuremath{\leftrightarrow}}{<->}file verification: declared versus recomputed row counts and checksum agreement per file, from verify\_descriptor\_files(). Both shipped fixture files verify against their bytes on disk; the title reports the descriptor's readiness score. If a fixture file were edited without updating the descriptor, the corresponding row would flip to a mismatch status and the test suite would fail.}]{../figures/checksum_verification.png}}
\caption{Descriptor\texorpdfstring{\ensuremath{\leftrightarrow}}{<->}file verification: declared versus recomputed row
counts and checksum agreement per file, from
\protect\breaktt{verify\_descriptor\_files()}. Both shipped fixture files
verify against their bytes on disk; the title reports the descriptor's
readiness score. If a fixture file were edited without updating the
descriptor, the corresponding row would flip to a mismatch status and
the test suite would fail.}\label{fig:checksum_verification}
\end{figure}

\subsection{Machine-readable release
manifest}\label{machine-readable-release-manifest}

\breaktt{build\_release\_bundle\_manifest()} packages the above into a
deterministic, JSON-ready manifest containing the schema fingerprint,
per-file checksums and row counts, field-level unit/bounds/enumeration
summaries, provenance steps, and the readiness verdict --- carrying
metadata and checksums, never the dataset bytes. The thin script
\href{../scripts/generate_release_artifacts.py}{\breaktt{scripts/generate\_release\_artifacts.py}}
writes this manifest (plus the readiness report and constraint summary)
under \texttt{output/reports/} for pre-publication review.

\subsection{Validation evidence}\label{validation-evidence}

Every capability above is exercised by the zero-mock \texttt{tests/}
suite: structural validation and negative controls, order-independent
schema fingerprints, byte-level verification against real temporary
files, figure-data preparation, and an integration test that runs the
figure script end to end and asserts real PNGs are written. Coverage
exceeds the 90\% project gate with no mocks.

\newpage

\section{Usage notes}\label{sec:usage}

\subsection{Regenerating the
artifacts}\label{regenerating-the-artifacts}

From the monorepo root, regenerate the figures and the release-review
artifacts:

\begin{Shaded}
\begin{Highlighting}[]
\ExtensionTok{uv}\NormalTok{ run python projects/templates/template\_data\_descriptor/scripts/generate\_figures.py}
\ExtensionTok{uv}\NormalTok{ run python projects/templates/template\_data\_descriptor/scripts/generate\_release\_artifacts.py}
\end{Highlighting}
\end{Shaded}

The first writes the five figures under \breaktt{manuscript/figures/};
the second writes the descriptor readiness report, field-constraint
summary, and metadata-only release manifest under
\texttt{output/reports/}. Both are thin orchestrators: all computation
lives in the tested \breaktt{src/data\_descriptor/} package.

\subsection{Forking in a real dataset}\label{forking-in-a-real-dataset}

To adapt this template to a real dataset:

\begin{enumerate}
\def\labelenumi{\arabic{enumi}.}
\tightlist
\item
  Replace the fixture files under \texttt{data/fixtures/} with your data
  files (or, for restricted data, leave them out and publish descriptor
  + checksums only).
\item
  Edit \breaktt{data/example\_descriptor.json}: update \texttt{name},
  \texttt{license}, the \texttt{files} inventory (path, media type,
  sha256 checksum, row count), the \texttt{fields} data dictionary,
  \texttt{primary\_key}, and the \texttt{provenance} chain.
\item
  Recompute each file's checksum and row count and record them in the
  descriptor. For present files, \breaktt{verify\_descriptor\_files()}
  will confirm the descriptor matches the bytes; a mismatch fails the
  test suite.
\item
  Extend field constraints (patterns, enumerations, numeric bounds,
  units) before publishing real data, and keep the provenance chain
  honest --- cite real tools, versions, and operators.
\end{enumerate}

Use \breaktt{scripts/audit/copy\_exemplar.py} from the monorepo to fork
the template cleanly, and keep \breaktt{domain\_profile.yaml} and
\breaktt{experiment\_plan.yaml} aligned with your dataset.

\subsection{Boundaries and claims}\label{boundaries-and-claims}

This is a template with a synthetic demonstration dataset. Its claims
are limited to dataset structure, schema completeness, provenance depth,
license state, byte-level integrity, and release readiness --- never to
any empirical effect in the data. When you fork it, keep that
discipline: a data descriptor earns trust by describing a dataset
precisely and verifiably, not by over-claiming what the dataset shows.

\newpage

\section{References}\label{sec:references}

This exemplar draws on established data-publishing practice: the FAIR
guiding principles for reusable data \citep{wilkinson2016fair},
datasheets-style structured dataset documentation
\citep{gebru2021datasheets}, the W3C PROV data model for provenance
\citep{moreau2013prov}, and the Frictionless Data Package container
format for machine-readable descriptors
\citep{frictionless2023datapackage}.

The bibliography lives in
\href{references.bib}{\breaktt{manuscript/references.bib}} and is read by
Pandoc during PDF render; every \texttt{{[}@key{]}} citation in the
manuscript is resolved against that file. To validate that
\texttt{references.bib} is syntactically clean and complete:

\begin{Shaded}
\begin{Highlighting}[]
\ExtensionTok{uv}\NormalTok{ run python }\AttributeTok{{-}m}\NormalTok{ infrastructure.reference.citation.cli validate }\DataTypeTok{\textbackslash{}}
\NormalTok{    projects/templates/template\_data\_descriptor/manuscript/references.bib }\AttributeTok{{-}{-}strict}
\end{Highlighting}
\end{Shaded}




\clearpage

\bibliography{references}
\end{document}
